\subsection{Циклы}

Циклы позволяют выполнять один и тот же блок кода многократно.

\subsubsection*{Цикл for}
Используется, когда известно количество повторений. Функция \texttt{range(N)} создает последовательность от 0 до N-1.

\subsubsection*{Цикл while}
Выполняется, пока условие истинно. Будьте осторожны: если условие никогда не станет ложным, получится <<бесконечный цикл>>.

\begin{codefile}[python]{python/examples/01_python_basics/06_loops.py}
\end{codefile}

\begin{itemize}
    \item \texttt{break} — немедленно прерывает цикл.
    \item \texttt{continue} — переходит к следующей итерации.
\end{itemize}

\begin{taskbox}[Обратный отсчет]
Напишите программу, которая считает от 10 до 1, а затем пишет <<Поехали!>>. Используйте цикл \texttt{for} или \texttt{while}.
\end{taskbox}

\subsubsection*{Задачи для самостоятельного решения}

\begin{enumerate}
    \item Выведите числа от 0 до 4 используя \texttt{range(5)}.
    \item Выведите слово "Привет" 3 раза.
    \item Выведите числа от 1 до 5 (используйте \texttt{range(1, 6)}).
    \item Выведите четные числа от 0 до 8 (\texttt{range(0, 9, 2)}).
    \item Найдите сумму чисел от 1 до 100.
    \item Выведите числа от 5 до 1 (\texttt{range(5, 0, -1)}).
    \item Выведите таблицу умножения на 5 (от 1*5 до 10*5).
    \item Дана строка "Дрон". Выведите каждую букву с новой строки.
    \item Используя \texttt{while}, выводите \texttt{x}, пока \texttt{x < 3}. Начальное \texttt{x = 0}.
    \item Ищите число 7 в диапазоне от 0 до 10. Если нашли — используйте \texttt{break}.
    \item Выведите числа от 0 до 4, пропустив 2 (используйте \texttt{continue}).
    \item Выведите квадраты чисел от 1 до 5.
    \item Посчитайте длину строки "Hello" с помощью цикла (без \texttt{len()}).
    \item Выведите 5 строк, в каждой по одной звездочке \texttt{*}.
    \item Выведите лесенку из звезд (1, 2, 3...).
\end{enumerate}
